\documentclass[11pt,a4paper]{article}

% Polskie znaki i kodowanie
\usepackage[utf8]{inputenc}
\usepackage[T1]{fontenc}
\usepackage[polish]{babel}
\usepackage{geometry}
\geometry{top=2.5cm, bottom=2.5cm, left=2.5cm, right=2.5cm}

% Pakiety graficzne, tabelaryczne i formatowania
\usepackage{hyperref}
\usepackage{booktabs}
\usepackage{tabularx}
\usepackage{enumitem}
\usepackage{xcolor}
\usepackage{titlesec}
\usepackage{listings}
\usepackage{amsmath}
\usepackage{fancyhdr}

% Konfiguracja nagłówków i stopek
\pagestyle{fancy}
\fancyhf{}
\fancyhead[L]{\small Projekt: Gra w Statki (C)}
\fancyhead[R]{\small Plan Działania i Architektura}
\fancyfoot[C]{\thepage}

% Styl linków i nagłówków
\hypersetup{
    colorlinks=true,
    linkcolor=blue!70!black,
    citecolor=blue!70!black,
    urlcolor=blue!70!black
}

\titleformat{\section}{\large\bfseries\color{blue!60!black}}{\thesection.}{0.8em}{}[\titlerule]
\titleformat{\subsection}{\normalsize\bfseries\color{black!85}}{\thesubsection.}{0.6em}{}
\titleformat{\subsubsection}{\small\bfseries\color{black!70}}{\thesubsubsection.}{0.5em}{}

% Konfiguracja listingu kodu
\lstset{
    basicstyle=\ttfamily\footnotesize,
    keywordstyle=\color{blue}\bfseries,
    commentstyle=\color{gray}\itshape,
    stringstyle=\color{teal},
    backgroundcolor=\color{black!3},
    frame=single,
    rulecolor=\color{black!15},
    breaklines=true,
    captionpos=b,
    showstringspaces=false
}

\title{
    \vspace{-1.5cm}
    \textbf{\huge Gra w Statki (Battleship) w języku C}\\[0.3cm]
    \large Kompleksowy Plan Działania, Architektura Sieciowa LAN oraz Harmonogram Wdrożenia
}
\author{\textbf{Dokumentacja Inżynierska Projektu}}
\date{\today}

\begin{document}

\maketitle
\thispagestyle{fancy}

\begin{abstract}
Niniejszy dokument przedstawia szczegółowy plan projektowy oraz architekturę oprogramowania dla aplikacji klasycznej gry w ,,Statki'' (Battleship), zrealizowanej w standardzie języka C (C11). Projekt uwzględnia modularną konstrukcję, wsparcie dla rozgrywki wieloosobowej w sieci lokalnej (LAN) oraz na dwóch terminalach (\texttt{localhost}), zaawansowany moduł sztucznej inteligencji (AI), pełną parametryzację planszy oraz floty statków, a także interfejs użytkownika z możliwością łatwej rozbudowy.
\end{abstract}

\tableofcontents
\vspace{0.5cm}
\hrule
\vspace{0.5cm}

\section{Wprowadzenie i Założenia Projektowe}
\subsection{Opis Koncepcji}
Celem projektu jest stworzenie elastycznej, modularnej gry strategicznej w języku C, która zachowuje tradycyjne reguły morskiej bitwy, wprowadzając nowoczesne elementy architektury oprogramowania: elastyczną konfigurację, komunikację międzyprocesową przez gniazda TCP oraz algorytmy sztucznej inteligencji.

\subsection{Główne Wymagania Funkcjonalne}
\begin{enumerate}[noitemsep]
    \item \textbf{Tryb Singleplayer (vs AI):} Rozgrywka z komputerem o zróżnicowanych strategiach (strzały losowe, algorytm \textit{Hunt \& Target}, wyszukiwanie oparte na parzystości pól).
    \item \textbf{Tryb Multiplayer (LAN / 2 terminale):}
    \begin{itemize}[noitemsep]
        \item Gra na dwóch niezależnych komputerach podłączonych do wspólnej sieci LAN.
        \item Gra między dwoma oknami terminala na tym samym komputerze przy wykorzystaniu interfejsu pętli zwrotnej (\texttt{127.0.0.1}).
    \end{itemize}
    \item \textbf{Menu Główne i GUI/TUI:} Interaktywne menu umożliwiające wybór trybu rozgrywki, dostosowanie parametrów oraz czytelny podgląd stanu własnej planszy i radaru trafień wroga.
    \item \textbf{Pełna Konfigurowalność Rozgrywki:}
    \begin{itemize}[noitemsep]
        \item Dynamiczny dobór wielkości planszy (np. od $6\times6$ do $20\times20$; domyślnie $10\times10$).
        \item Możliwość definiowania liczebności floty oraz rozmiaru poszczególnych jednostek (np. jednomasztowce, dwumasztowce, trójmasztowce, czteromasztowce).
    \end{itemize}
    \item \textbf{Elastyczność Architektury:} Możliwość wprowadzania zmian i rozszerzeń w toku projektu (np. zmiana mechanizmu wyświetlania z TUI ANSI na bibliotekę graficzną Raylib/SDL, dodanie zapisu stanu gry do pliku).
\end{enumerate}

\section{Architektura Systemu i Podział Modularny}
Aby spełnić wymóg adaptacyjności kodu, zastosowano architekturę warstwową inspirowaną wzorcem MVC (Model-View-Controller) z wyraźnym odseparowaniem logiki biznesowej od operacji wejścia/wyjścia oraz warstwy transportowej.

\begin{table}[h!]
\centering
\small
\begin{tabularx}{\textwidth}{l p{4.0cm} X}
\toprule
\textbf{Moduł} & \textbf{Pliki źródłowe} & \textbf{Zakres odpowiedzialności} \\
\midrule
\textbf{Config} & \texttt{include/config.h} & Globalne stałe, limity, domyślne parametry planszy i floty statków. \\
\addlinespace[2pt]
\textbf{Game Logic} & \texttt{include/game.h} \newline \texttt{src/game.c} & Dynamiczne tworzenie i niszczenie planszy, walidacja reguł rozstawiania statków, rejestracja strzałów, sprawdzanie warunku zwycięstwa. \\
\addlinespace[2pt]
\textbf{Network} & \texttt{include/network.h} \newline \texttt{src/network.c} & Abstrakcja gniazd TCP (cross-platform: Windows Winsock2 / Linux POSIX Sockets), nawiązywanie połączeń, wysyłanie i odbiór pakietów. \\
\addlinespace[2pt]
\textbf{Protocol} & \texttt{include/protocol.h} \newline \texttt{src/protocol.c} & Serializacja i deserializacja komunikatów sieciowych (uzgadnianie parametrów gry, koordynaty strzałów, wyniki). \\
\addlinespace[2pt]
\textbf{AI} & \texttt{include/ai.h} \newline \texttt{src/ai.c} & Procedury automatycznego rozstawiania statków oraz silnik generowania ruchów bota (Hunt \& Target). \\
\addlinespace[2pt]
\textbf{UI / TUI} & \texttt{include/ui.h} \newline \texttt{src/ui.c} & Prezentacja menu, renderowanie plansz obok siebie z kolorystyką ANSI, obsługa wprowadzania współrzędnych. \\
\addlinespace[2pt]
\textbf{Controller} & \texttt{src/main.c} & Pętla główna, maszyna stanów gry, orkiestracja przepływu danych pomiędzy modułami. \\
\bottomrule
\end{tabularx}
\caption{Podział modularny projektu gry w Statki}
\end{table}

\section{Kluczowe Struktury Danych}
Wszystkie tablice planszy alokowane są dynamicznie w stercie, co zapewnia brak sztywnych ograniczeń kompilacyjnych:

\begin{lstlisting}[language=C, caption={Podstawowe struktury silnika gry (game.h)}]
typedef enum {
    CELL_EMPTY = 0, // Woda / puste pole
    CELL_SHIP,      // Segment statku (ukryty przed wrogiem)
    CELL_MISS,      // Pudlo (strzal w wode)
    CELL_HIT,       // Trafiony segment statku
    CELL_SUNK       // Segment zatopionego statku
} CellState;

typedef struct {
    int id;             // Identyfikator statku
    int x, y;           // Pozycja dziobu
    int size;           // Dlugosc jednostki
    Direction dir;      // DIR_HORIZONTAL lub DIR_VERTICAL
    int hits_left;      // Licznik nieuszkodzonych pol
    bool is_sunk;       // Flaga zatopienia
} Ship;

typedef struct {
    int width, height;  // Dynamiczne wymiary planszy
    CellState **grid;   // Dynamiczna macierz pol [height][width]
    int **ship_ids;     // Odwzorowanie pola na ID statku
    Ship *ships;        // Tablica statkow gracza
    int ship_count;     // Calkowita liczba statkow
    int ships_alive;    // Liczba statkow plywajacych
} Board;
\end{lstlisting}

\section{Komunikacja Sieciowa i Protokół LAN}
\subsection{Topologia Połączenia}
Rozgrywka sieciowa wykorzystuje protokół \textbf{TCP/IP}. 
\begin{itemize}[noitemsep]
    \item \textbf{Gracz 1 (Host):} Inicjuje gniazdo serwerowe (\texttt{net\_listen}) na ustalonym porcie (domyślnie \texttt{5555}) i oczekuje na połączenie.
    \item \textbf{Gracz 2 (Klient):} Nawiązuje połączenie (\texttt{net\_connect}), podając adres IP serwera (np. \texttt{192.168.1.15} dla sieci LAN lub \texttt{127.0.0.1} dla gry na dwóch terminalach tego samego komputera).
\end{itemize}

\subsection{Przebieg Wymiany Danych (Handshake i Tury)}
Protokół wymiany opiera się na czytelnych ramkach tekstowych zakończonych znakiem nowej linii \texttt{\textbackslash n}:

\begin{enumerate}
    \item \textbf{Faza Negocjacji (Handshake):}
    \begin{itemize}[noitemsep]
        \item \texttt{Host $\rightarrow$ Klient: CONFIG <szerokość> <wysokość> <liczba\_statków> <s1> <s2> ...}
        \item \texttt{Klient $\rightarrow$ Host: CONFIG\_ACK}
    \end{itemize}
    \item \textbf{Faza Gotowości:}
    Obaj gracze rozstawiają statki (ręcznie lub automatycznie). Po zakończeniu każdy wysyła:
    \begin{itemize}[noitemsep]
        \item \texttt{Gracz $\rightarrow$ Rywal: READY}
    \end{itemize}
    \item \textbf{Faza Rozgrywki (Naprzemienne tury):}
    \begin{itemize}[noitemsep]
        \item \texttt{Gracz A $\rightarrow$ Gracz B: SHOT <x> <y>}
        \item Gracz B weryfikuje strzał na swojej planszy i odpowiada:
        \item \texttt{Gracz B $\rightarrow$ Gracz A: RESULT <x> <y> <KOD>} (gdzie \texttt{KOD} to: \texttt{MISS}, \texttt{HIT}, \texttt{SUNK} lub \texttt{WIN}).
        \item Jeżeli nastąpiło pudło, tura przechodzi na drugiego gracza; przy trafieniu gracz kontynuuje ruch.
    \end{itemize}
    \item \textbf{Zakończenie Gry:}
    Komunikat \texttt{WIN} lub rezygnacja gracza (\texttt{QUIT}).
\end{enumerate}

\section{Algorytm Sztucznej Inteligencji (AI)}
Dla trybu jednoosobowego zaimplementowano algorytm działający w dwóch stanach:
\begin{enumerate}
    \item \textbf{Faza Poszukiwania (Hunt Mode):} Bot oddaje strzały w pola o parzystej sumie współrzędnych $(x + y) \pmod 2 = 0$ (heurystyka szachownicy). Ponieważ najmniejszy wielomasztowiec ma rozmiar $\ge 2$, ta taktyka pozwala na dwukrotne przyspieszenie wykrycia pierwszego segmentu statku.
    \item \textbf{Faza Niszczenia (Target Mode):} Po zarejestrowaniu trafienia (\texttt{HIT}) bot wstrzymuje losowe strzały i bada 4 sąsiadujące pola (góra, prawo, dół, lewo). Po kolejnym trafieniu ustala oś statku (pozioma/pionowa) i kontynuuje ostrzał wzdłuż tej osi aż do otrzymania komunikatu \texttt{SUNK}.
\end{enumerate}

\section{Harmonogram Realizacji (Kroki Implementacji)}
Projekt został rozłożony na pięć powiązanych etapów:

\subsection*{Etap 1: Silnik Podstawowy i Mechanika Gry}
\begin{itemize}[noitemsep]
    \item Implementacja dynamicznej alokacji i zwalniania struktur \texttt{Board} i \texttt{Ship}.
    \item Algorytm walidacji pozycji statku (brak kolizji i brak stykania się rogami ani burtami).
    \item Mechanizm obsługi strzału, aktualizacji stanu segmentów i wykrywania zatopienia całej floty.
\end{itemize}

\subsection*{Etap 2: Interfejs Użytkownika (TUI) i Formularz Konfiguracji}
\begin{itemize}[noitemsep]
    \item Ekran Menu Głównego oraz podmenu Ustawień (konfiguracja rozmiaru planszy i floty).
    \item Rysowanie podwójnej planszy w konsoli (plansza własna ze statkami + plansza trafień).
    \item Parsowanie wejścia użytkownika (obsługa notacji szachowej np. \texttt{A5} oraz liczbowej \texttt{1 5}).
\end{itemize}

\subsection*{Etap 3: Logika Sztucznej Inteligencji (AI)}
\begin{itemize}[noitemsep]
    \item Generator automatycznego, bezkolizyjnego rozstawiania losowej floty statków.
    \item Implementacja maszyny stanów bota (\textit{Hunt \& Target}).
    \item Pętla rozgrywki w trybie \textit{Singleplayer} przeciwko komputerowi.
\end{itemize}

\subsection*{Etap 4: Moduł Sieciowy i Tryb Wieloosobowy (LAN / 2 Terminale)}
\begin{itemize}[noitemsep]
    \item Implementacja warstwy gniazd TCP z obsługą wieloplatformowości (Winsock / POSIX).
    \item Wdrożenie serializacji i deserializacji protokołu sieciowego.
    \item Integracja maszyny stanów dla ról Hosta i Klienta, obsługa timeoutów i rozłączeń.
    \item Testy rozgrywki między dwoma oknami terminala oraz w fizycznej sieci LAN.
\end{itemize}

\subsection*{Etap 5: Rozszerzenia, Obsługa Zmian i Ostateczny Szlif}
\begin{itemize}[noitemsep]
    \item Możliwość podpięcia alternatywnego widoku (np. okienkowego GUI opartego o Raylib lub SDL).
    \item Zabezpieczenie przed wyciekami pamięci (weryfikacja narzędziami Valgrind / Dr. Memory).
    \item Testy scenariuszy brzegowych (błędne pakiety, zerwanie łącza).
\end{itemize}

\section{Zarządzanie Zmianami w Projekcie}
W toku realizacji wymagania mogą ulegać zmianie. Architektura została przygotowana na następujące modyfikacje:
\begin{itemize}[noitemsep]
    \item \textbf{Zmiana biblioteki interfejsu:} Dzięki odseparowaniu \texttt{ui.h} od \texttt{game.h}, zastąpienie terminala interfejsem graficznym nie wymaga żadnych modyfikacji w silniku gry ani w module sieciowym.
    \item \textbf{Nowe typy statków / specjalne zdolności:} Struktura \texttt{Ship} zawiera identyfikator oraz elastyczne pole długości, co pozwala łatwo dodać np. miny morskie czy samoloty zwiadowcze.
    \item \textbf{Rozbudowa protokołu sieciowego:} Protokół tekstowy umożliwia bezproblemowe dołożenie nowych komend, np. \texttt{CHAT <tekst>} czy \texttt{PAUSE}.
\end{itemize}

\section{Podsumowanie i Instrukcja Kompilacji}
Projekt posiada gotowe konfiguracje dla dwóch najpopularniejszych narzędzi budowania:
\begin{itemize}
    \item \textbf{CMake:}
    \begin{lstlisting}[language=bash]
mkdir build && cd build
cmake ..
cmake --build .
    \end{lstlisting}
    \item \textbf{GNU Make (GCC / MinGW):}
    \begin{lstlisting}[language=bash]
make
./statki.exe  # (lub ./statki na systemach Linux/macOS)
    \end{lstlisting}
\end{itemize}

\end{document}
